% Os apêndices são materiais elaborados pelo autor que complementam sua argumentação,
% mas que não são essenciais para a compreensão do texto principal.
% Exemplos: demonstrações matemáticas extensas, códigos-fonte, especificações
% detalhadas, documentação técnica, etc.

\begin{apendicesenv}

\partapendices

\chapter{Pesquisa Bibliográfica}
\label{ap-pesquisa-bibliografica}

Este apêndice detalha o levantamento bibliográfico realizado para fundamentar esta pesquisa.
Para cada obra, apresenta-se uma análise crítica que abrange o
contexto de aplicação, as contribuições metodológicas ou tecnológicas pertinentes
a este trabalho, além das limitações ou restrições de escopo que caracterizam as
lacunas a serem exploradas por esta monografia. A \autoref{tab:matriz-comparativa},
ao final do apêndice, sintetiza as dimensões avaliadas em formato comparativo.

% ============================================================
\section{SIGDENGUE}
\label{ap:sigdengue-linha}
% ============================================================

Os dois artigos a seguir apresentam ligação direta e correspondem a uma linha de
pesquisa iniciada em 2011 na Universidade Estadual do Oeste do Paraná (UNIOESTE),
em parceria com a Prefeitura de Cascavel/PR. Estes representam os estudos mais
alinhados aos objetivos desta pesquisa.

% ============================================================
\subsection{\textit{Computational Agents Applied to Dengue Fever Simulation} (2011)}
\label{ap:rizzi2011}
% ============================================================

\citeonline{rizzi2011} apresentam o estágio de desenvolvimento de um
sistema de monitoramento de dengue composto por três módulos, que visam:

\begin{enumerate}[label=\Roman*.]
    \item \textbf{Sistema de Informação:} Integrar dados e informações sobre
    ocorrências de dengue de forma ordenada e centralizada, além de agregar
    informações meteorológicas, densidade populacional e distribuição geográfica
    no município;
    \item \textbf{Georreferenciamento:} Fornecer ferramentas que facilitem a
    realização de simulações da disseminação da dengue. Utiliza mapas no formato
    \textit{ESRI Shapefile}, contendo três arquivos: desenho do mapa (.shp), índice
    de referência para o arquivo .shp e atributos do mapa (.dbf);
    \item \textbf{Simulação com Agentes:} Modelagem do ambiente e dos agentes
    (mosquitos e humanos).
\end{enumerate}

O \textbf{Módulo~III} é o ponto de maior relevância do artigo, pois discorre sobre a
simulação que emprega um modelo compartimental
(\textbf{SEIR}\footnote{A população é dividida em quatro compartimentos:
\textbf{S}uscetível (indivíduos saudáveis e expostos ao risco), \textbf{E}xposto
(infectados, mas ainda não transmissores, no período de incubação), \textbf{I}nfeccioso
(transmissores ativos da doença) e \textbf{R}ecuperado (imunes temporária ou
permanentemente). O fluxo entre compartimentos é definido por taxas paramétricas.
Para a dengue, a extensão habitual é o modelo SEIRS, que inclui perda de imunidade
e retorno ao estado suscetível.}) sobre a população, buscando representar as
conexões sociais reais dos indivíduos (casa, trabalho, estudo e lazer) e seus
deslocamentos físicos.

Os mosquitos possuem sensor visual para a percepção de corpos no ambiente, além
de sensores de temperatura e odor para captar as condições que afetam sua
capacidade de sobrevivência. Os agentes foram desenvolvidos seguindo a arquitetura
\textbf{BDI} (\textit{Belief-Desire-Intention}) e os protocolos de mensagem oferecidos
pelo \textbf{JADE}. Os autores afirmam
que o baixo nível de cognição decorre da abrangência do problema, mas que os padrões
de comunicação e relacionamento são considerados essenciais, com possibilidade de
melhoria na modelagem. Esse aspecto constitui o principal gancho aproveitado neste
trabalho.

O primeiro módulo é explorado mais a fundo no artigo subsequente
\cite{rizzi2016}, enquanto o segundo módulo não recebe o mesmo nível de
aprofundamento.

% ============================================================
\subsection{SIGDENGUE: Um Sistema de Informação para o Acompanhamento e Gestão
de Ações sobre Dengue com Enfoque às Atividades de Notificação, Raio e Bloqueio
(2016)}
\label{ap:rizzi2016}
% ============================================================

\citeonline{rizzi2016} descrevem o \textit{Sistema de Informação para
Aquisição, Manipulação e Tratamento de Dados sobre a Dengue} (\textbf{SIGDENGUE}),
desenvolvido para integrar as bases oficiais brasileiras de vigilância da dengue
(SINAN, LIRAa e SISPNCD) em uma única plataforma municipal com suporte a
georreferenciamento e tomada de decisão. O sistema visa facilitar o acesso à
informação por meio de consultas, relatórios e visualização georreferenciada das
ocorrências, consolidando dados que o município parceiro dispunha de forma
fragmentada e os complementando com os registros dos sistemas nacionais. A
arquitetura foi concebida de modo a ser adaptável a outros municípios.

Lido em conjunto com seu antecessor \cite{rizzi2011}, o par de artigos
oferece uma visão complementar da linha de pesquisa: enquanto o trabalho de 2011
explora a simulação com agentes e expõe lacunas na profundidade da modelagem
comportamental, o de 2016 amplia o escopo para o sistema de informação e oferece
um panorama abrangente sobre os dados disponíveis na vigilância da dengue no Brasil.
A revisão bibliográfica apresentada neste segundo artigo foi um recurso valioso para
mapear as fontes de dados existentes e compreender como podem ser utilizadas como
insumos para a simulação proposta neste trabalho.

% ============================================================
\section{Modelagem da Dengue}
\label{ap:abm-mosquitos}
% ============================================================

Os artigos agrupados nesta seção destacam-se pelo enfoque na modelagem de
arboviroses. A análise desses trabalhos é fundamental para a compreensão das regras
biológicas e das dinâmicas de transmissão, pois apresentam diferentes abstrações e
estratégias de simulação que servem de referência para a modelagem dos agentes deste
simulador.

\subsection{\textit{Exploring the Dynamics of Gene Drive Mosquitoes Within Wild
Populations Using an Agent-Based Simulation} (2024)}
\label{ap:wickramasooriya2024}

\citeonline{wickramasooriya2024} propõem um modelo ABM geoespacial para simular
a liberação de mosquitos geneticamente modificados (\textit{gene drive}) na ilha do
Príncipe, África Ocidental, com o objetivo de controlar a malária transmitida pelo
\textit{Anopheles coluzzii}. O modelo é organizado em três camadas:

\begin{enumerate}[label=(\alph*)]
    \item \textbf{Ciclo de vida:} transições de estado entre ovo, larva, pupa e adulto,
    com probabilidades de sobrevivência ajustadas por dados laboratoriais obtidos do
    laboratório de genética vetorial da Universidade da Califórnia, Davis;
    \item \textbf{Movimentação e dispersão:} voo baseado em custo de resistência
    calculado a partir da elevação do terreno e do tipo de uso do solo (dados GIS),
    com alcance diário entre 170~m e 1.736~m. O custo mínimo de resistência é otimizado a cada deslocamento para reproduzir a preferência
    ecológica do mosquito por regiões de baixa altitude e determinados tipos de
    cobertura;
    \item \textbf{Acasalamento e oviposição:} comunicação por troca de mensagens
    entre agentes para seleção de parceiros e transmissão de genótipo. A fêmea
    enfileira pedidos de machos disponíveis em seu raio de 50~m e seleciona
    aleatoriamente, transitando para os estados de busca de sangue e oviposição
    na sequência.
\end{enumerate}

Os diagramas de estado para fêmeas e machos são formalmente especificados,
e o trabalho relata testes unitários, de integração e de sistema como práticas
explícitas de engenharia de \textit{software}. A separação em três camadas
especializadas com diagramas de estado formais e o algoritmo de movimentação
baseado em resistência geoespacial são diretamente adaptáveis ao
\textit{Aedes aegypti}. A estratégia de testes em três níveis também serve como
uma excelente referência de verificação para este projeto. Entretanto, os próprios autores apontam a ausência de variáveis climáticas
(temperatura, chuva, umidade) no modelo de sobrevivência, fator crítico para o
\textit{Aedes aegypti}.

% ============================================================
\subsection{\textit{Agent-Based Modeling and Simulation for Transmission Dynamics
and Surveillance of Dengue: Conceptual and Design Model} (2023)}
\label{ap:pascoe2023}
% ============================================================

\citeonline{pascoe2023} apresentam um modelo conceitual e de \textit{design}
para a vigilância da dengue na região de Dar es Salaam (Tanzânia), utilizando a
plataforma MARS (\textit{Multi-Agent Research and Simulation}). O modelo segue o
protocolo ODD (\textit{Overview, Design Concepts, and Details}), padrão da
comunidade de modelagem baseada em agentes para garantir a reprodutibilidade.

Os agentes são a fêmea do \textit{Aedes aegypti} e o humano. O ambiente é
representado por camadas GIS em formato \textit{raster} cobrindo temperatura,
precipitação e umidade. A movimentação dos mosquitos segue o
\textit{Movement Ecology Paradigm} (MEP); a dos humanos incorpora planos diários
(residência, escola, trabalho e lazer). O modelo SEIR é adotado para a progressão
da doença no agente humano, com parâmetros extraídos da literatura (período de
incubação de 4 a 10 dias, taxa de picada do mosquito de 0,26 a 0,67/dia, vida útil
do mosquito de 30 dias).

A parametrização explícita de todas as variáveis de estado dos agentes, as
modalidades de movimentação e a aplicação rigorosa do protocolo ODD são referências
diretas para a especificação dos agentes deste simulador. Contudo, o trabalho
permanece no nível de projeto. Não há código implementado, resultados de simulação
ou validação com dados empíricos reais.

% ============================================================
% TODO: Revisar o artigo
\subsection{\textit{Mixed-mode pandemic modeling} (2019)}
\label{ap:bosch2019}

\citeonline{bosch2019} descrevem uma arquitetura de simulação híbrida para
pandemias globais que combina três paradigmas: Dinâmica de Sistemas (SD), Simulação
de Eventos Discretos (DES) e Modelagem Baseada em Agentes (ABM). Cada célula de
30~km² da superfície terrestre contém um modelo SD para a progressão da doença;
um motor DES coordena a evolução temporal e os eventos de viagem aérea; e agentes
ABM representam governos e equipes de resposta, tomando decisões sobre vacinação,
quarentena e alocação de recursos. Alguns pontos relevantes podem ser destacados, como:

\begin{enumerate}
    \item \textbf{Distinção formal entre endemia e epidemia:} o modelo define estados
    endêmicos como aqueles em que a taxa de incidência ocorre de forma estável e
    esperada, e estados epidêmicos como aqueles em que essa taxa supera o limite do
    esperado. Essa distinção fundamenta a definição de surto adotada neste simulador
    e os critérios de disparo de alertas para os gestores de saúde pública;
    \item \textbf{Argumento contra modelos puramente compartimentais:} o trabalho
    demonstra que os modelos SD e SIR/SEIR, por operarem com populações médias e
    homogêneas, são insuficientes para capturar as decisões individuais de agentes
    autônomos (governos, equipes de campo). Esse argumento reforça diretamente a
    escolha pelo paradigma ABM nesta pesquisa.
\end{enumerate}

A limitação identificada é que o modelo é declarado como uma prova de conceito sem
validação com dados reais. Além disso, a granularidade espacial de 30~km² é
inadequada para simulações intraurbanas, onde criadouros e focos de dengue exigem
uma resolução na escala de dezenas de metros.

% ============================================================
\section{Dados Epidemiológicos e Ciência de Dados}
\label{ap:dados}
% ============================================================

Esta seção agrupa estudos que aplicam técnicas de Ciência de Dados e mineração
estatística ao tratamento de bases epidemiológicas. A análise desses trabalhos
orienta a seleção, normalização e validação dos dados reais utilizados como
insumos do simulador, garantindo uma parametrização baseada em evidências.

% ============================================================
% TODO: Revisar o artigo
\subsection{\textit{Risk Factor Identification and Spatiotemporal Diffusion Path
During the Dengue Outbreak} (2016)}
\label{ap:li2016}
% ============================================================

\citeonline{li2016} analisam a epidemia de dengue ocorrida no Distrito de Xiangqiao
(Chaozhou, China) em 2015, utilizando uma abordagem de mineração de dados
espaço-temporais em nível microscópico (resolução temporal diária). O trabalho foca
na extração de conhecimento a partir de grandes volumes de dados para identificar
variáveis críticas. A metodologia combina:

\begin{enumerate}
    \item \textbf{Análise Estatística de Ciclo:} Ajuste de curvas por mínimos
    quadrados para definir o ciclo de vida epidemiológico em quatro fases
    (\textit{initial}, \textit{beginning}, \textit{outbreak} e \textit{extinction});
    \item \textbf{Autocorrelação Espacial:} Uso do índice de Moran para quantificar
    a dependência espacial dos casos, validando a necessidade de representação
    geográfica;
    \item \textbf{Correlação de Fatores:} Aplicação da correlação de Spearman para
    medir a associação entre a incidência de dengue e fatores de risco (população,
    temperatura, umidade, vento e Índice de Breteau);
    \item \textbf{Geoprocessamento:} Mapeamento de densidade de Kernel para
    identificar \textit{hotspots} e trajetórias de difusão.
\end{enumerate}

Os resultados derivados da análise de dados indicam que o Índice de Breteau (IB) e
a densidade populacional são os fatores de risco dominantes durante o surto ativo.
Notavelmente, as variáveis meteorológicas diárias apresentaram correlação fraca no
curto prazo, resultado contraintuitivo que orienta a priorização de variáveis de
entrada no simulador. Além disso, o Índice de Moran superior a 0,52 confirma uma
forte autocorrelação espacial, fornecendo evidências de que a vizinhança geográfica
deve ser tratada como uma variável de estado explícita no SMA.

% ============================================================
\subsection{\textit{Data Science Approaches for Building Dengue Epidemiological
Datasets} (2025)}
\label{ap:moreira2025}
% ============================================================

\citeonline{moreira2025} constroem e disponibilizam \textit{datasets}
normalizados de dengue para três cidades brasileiras (São José dos Campos,
Ouro Preto e Resende), no período de 2014 a 2024, a partir de dados
do SINAN e do InfoDengue. Os \textit{datasets} normalizados estão disponíveis
publicamente no repositório Zenodo em formato CSV, acessíveis para uso direto
como dados de validação.

Para o SINAN, são propostas três relações normalizadas cobrindo casos
autóctones\footnote{Os casos autóctones referem-se a indivíduos que contraem a
doença em seu local de residência, sem histórico de viagens para áreas com
transmissão conhecida. Esta métrica serve como um indicador fundamental para
monitorar a propagação local da doença.}, hospitalizações e desfechos clínicos.
Na normalização das bases para o InfoDengue, a relação construída foi:
$$R_1 = \{\text{epi\_weeks}, \text{cases}, \text{avg\_temp}, \text{avg\_humid}\}$$

A Função de Correlação Cruzada (CCF) entre temperatura e incidência apresentada
revela um \textbf{defasamento de aproximadamente 12 semanas} (três meses) entre o
pico de temperatura e o pico de casos, dado de suma importância para a
parametrização dos agentes mosquitos neste simulador. Além disso, foi confirmada a
ocorrência de ciclos endêmicos a cada 4 a 5 anos, evidenciando um padrão que deve
ser reprodutível pelo simulador como critério de validação temporal. Cabe destacar
que a ausência de dados pluviométricos (chuva) no InfoDengue é apontada pelos
próprios autores como uma lacuna relevante.

% ============================================================
\section{Matriz Comparativa}
\label{ap:matriz}
% ============================================================
% TODO: Revisar tabela

A \autoref{tab:matriz-comparativa} consolida a análise dos trabalhos abordados em
seis dimensões consideradas para os objetivos deste TCC. A marcação ``X''
indica que a dimensão está contemplada no trabalho (seja de forma total ou parcial),
enquanto os espaços em branco indicam a ausência da característica.

\begin{table}[htb]
    \centering
    \caption{Matriz comparativa dos trabalhos relacionados.}
    \label{tab:matriz-comparativa}
    \vspace{0.2cm}
    \begin{tabular}{|p{4.cm}|c|c|c|c|c|c|c|}
        \hline
        \multirow{2}{4.5cm}{\textbf{Dimensões Analisadas}} &
        \multicolumn{7}{c|}{\textbf{Trabalhos Relacionados (Seções)}} \\
        \cline{2-8}
        & \textbf{A.1.1} & \textbf{A.1.2} & \textbf{A.2.1} & \textbf{A.2.2} &
        \textbf{A.2.3} & \textbf{A.2.4} & \textbf{A.3.1} \\
        \hline
        Modelagem Baseada em Agentes (ABM) & X &   & X & X & X &   &   \\
        \hline
        Integração com Dados Abertos       & X & X &   &   &   &   & X \\
        \hline
        Uso de Variáveis Climáticas        & X &   &   & X &   & X & X \\
        \hline
        Georreferenciamento (GIS)          & X & X & X & X & X & X &   \\
        \hline
        Validação com Dados Reais          &   & X &   &   &   & X &   \\
        \hline
        Foco no Contexto Brasileiro        & X & X &   &   &   &   & X \\
        \hline
    \end{tabular}
    \par\vspace{1ex}
    {\footnotesize \textbf{Fonte:} Autor.}
\end{table}

A matriz evidencia que nenhum dos trabalhos isolados identificados cobre as seis
dimensões simultaneamente. A proposta desta pesquisa busca
integrá-las de forma conjunta, o que justifica a sua contribuição
científica para o estado da arte.

\end{apendicesenv}